\chapter{Simetri}\label{ch:g4:symmetry}

Kupu-kupu, wajah manusia, kepingan salju: alam menyukai bangun yang
kedua paruhnya berimpit ketika dilipat. Bab ini menjelajahi simetri
dengan lipatan dan kertas berpetak --- lukisan dengan penggaris siku
menunggu pada \cref{ch:g6:symmetry}.

\section{Melipat dan mencocokkan}

\begin{definition}[Sumbu simetri]\label{def:g4:symmetry:def}
Sebuah garis disebut \emph{sumbu simetri}\index{sumbu simetri}
sebuah bangun bila melipat bangun itu sepanjang garis tersebut
membuat kedua paruhnya berimpit \emph{tepat}. Sebuah bangun dapat
tidak punya sumbu, punya satu sumbu, atau punya beberapa sumbu.
\end{definition}

\begin{omfigure}
\begin{tikzpicture}[scale=0.6]
  % bangun simetris mirip hati: segitiga sama kaki
  \draw[thick, fill=omDef!12] (0,0) -- (2.4,0) -- (1.2,2.6) -- cycle;
  \draw[omThm, dashed, thick] (1.2,-0.5) -- (1.2,3.0);
  \node[below, font=\small] at (1.2,-0.7) {satu sumbu};
  % persegi panjang: dua sumbu
  \begin{scope}[xshift=5.4cm]
  \draw[thick, fill=omDef!12] (0,0.3) rectangle (3.2,2.3);
  \draw[omThm, dashed, thick] (1.6,-0.4) -- (1.6,3.0);
  \draw[omThm, dashed, thick] (-0.5,1.3) -- (3.7,1.3);
  \node[below, font=\small] at (1.6,-0.7) {dua sumbu};
  \end{scope}
  % segitiga sembarang: tanpa sumbu
  \begin{scope}[xshift=11.2cm]
  \draw[thick, fill=omDef!12] (0,0) -- (3.0,0.4) -- (0.8,2.2) -- cycle;
  \node[below, font=\small] at (1.4,-0.7) {tanpa sumbu};
  \end{scope}
\end{tikzpicture}

{\small Lipatlah sepanjang garis putus-putus: kedua paruhnya
berimpit. Segitiga terakhir tidak punya lipatan yang berhasil --- ia tanpa \omterm{def:g4:symmetry:def}{sumbu simetri}.}
\end{omfigure}

\begin{example}[Menguji dengan kertas kalkir]\label{ex:g4:symmetry:trace}
Jiplaklah gambarnya, baliklah kertas kalkirnya sepanjang sumbu yang
diduga, lalu letakkan kembali: kalau jiplakannya menutupi gambar itu
dengan tepat, sumbunya asli. Mata saja sering tertipu --- diagonal
sebuah persegi panjang \emph{terlihat} seperti sumbu, tetapi uji
lipatan mengatakan tidak (cobalah!).
\end{example}

\section{Melengkapi gambar di kertas berpetak}

\begin{method}[Bayangan sebuah gambar terhadap garis petak]\label{met:g4:symmetry:grid}
Sumbunya adalah garis tegak (atau mendatar) pada petak itu.
\begin{enumerate}
  \item Ambillah setiap \omterm{def:g4:geometry:polygon}{titik sudut} gambar itu satu per satu;
  \item bilanglah jaraknya ke sumbu dalam satuan petak;
  \item letakkan \omterm{def:g4:geometry:polygon}{titik sudut} bayangannya pada jarak yang \emph{sama}
        di sisi yang \emph{lain}, pada baris (atau kolom) yang sama;
  \item hubungkan \omterm{def:g4:geometry:polygon}{titik sudut} bayangannya dengan urutan yang sama,
        lalu periksa: bayangannya adalah kembaran cermin gambar itu.
\end{enumerate}
\end{method}

\begin{omfigure}
\begin{tikzpicture}[scale=0.55]
  \draw[gray!40, thin] (0,0) grid (12,5);
  \draw[omThm, very thick] (6,-0.4) -- (6,5.4);
  \draw[thick, omDef, fill=omDef!15]
    (1,1) -- (5,1) -- (5,2) -- (3,2) -- (3,4) -- (1,4) -- cycle;
  \draw[thick, omMeth, fill=omMeth!15]
    (11,1) -- (7,1) -- (7,2) -- (9,2) -- (9,4) -- (11,4) -- cycle;
  \draw[dashed] (5,2) -- (7,2);
  \draw[dashed] (3,4) -- (9,4);
\end{tikzpicture}

{\small Melengkapi terhadap sumbu merah: setiap \omterm{def:g4:geometry:polygon}{titik sudut} melompat
ke jarak yang sama di sisi yang lain. \omterm{def:g4:geometry:polygon}{Titik sudut} yang menyentuh
sumbu akan tetap di tempatnya.}
\end{omfigure}

\begin{example}[Bentuk sama, ukuran sama, tetapi terbalik]\label{ex:g4:symmetry:properties}
Bayangannya punya panjang dan sudut yang persis sama dengan aslinya:
simetri menyalin gambar --- tetapi membaliknya, seperti tangan kiri
dan tangan kanan. Kalau huruf aslinya F, bayangan cerminnya adalah F
yang terbalik, bukan F yang lain.
\end{example}

\section{Sumbu bangun-bangun yang biasa}

\begin{example}[Membilang sumbu]\label{ex:g4:symmetry:shapes}
Dengan melipat:
\begin{itemize}
  \item persegi: $4$ sumbu (dua melalui tengah sisi yang berhadapan,
        dua sepanjang diagonalnya);
  \item persegi panjang: $2$ sumbu (hanya melalui tengah sisi yang
        berhadapan!);
  \item segitiga sama kaki: $1$; segitiga sama sisi: $3$;
  \item lingkaran: setiap garis yang melalui pusatnya --- sumbunya
        lebih banyak daripada yang dapat dibilang.
\end{itemize}
\end{example}

\begin{example}[Simetri pada abjad]\label{ex:g4:symmetry:letters}
A, M, T, U, V punya sumbu tegak; B, C, D, E punya sumbu mendatar;
H, I, O, X punya keduanya. Kata yang tersusun dari huruf yang tepat
juga bisa simetris: MUM punya sumbu tegak; OXO berhasil ke dua arah.
\end{example}

\section{Latihan}

\begin{exercise}[$\star$]\label{exo:g4:symmetry:1}
Jiplak lalu lipatlah: mana di antara bangun ini yang punya \omterm{def:g4:symmetry:def}{sumbu
simetri} --- persegi; persegi panjang (yang bukan persegi); segitiga
sembarang (semua sisinya berbeda); lingkaran?
\end{exercise}

\begin{exercise}[$\star$]\label{exo:g4:symmetry:2}
Gambarlah \omterm{def:g4:symmetry:def}{sumbu simetri} dari: persegi; persegi panjang; segitiga
sama sisi. Berapa sumbu untuk masing-masing?
\end{exercise}

\begin{exercise}[$\star$]\label{exo:g4:symmetry:3}
Pakailah uji kertas kalkir pada \cref{ex:g4:symmetry:trace} untuk
menunjukkan bahwa diagonal persegi panjang $6 \times 4$ \emph{bukan}
\omterm{def:g4:symmetry:def}{sumbu simetri}.
\end{exercise}

\begin{exercise}[$\star$]\label{exo:g4:symmetry:4}
Di kertas berpetak, gambarlah sumbu tegak dan huruf L (setinggi tiga
petak, selebar dua petak) pada jarak $2$ petak dari sumbu itu.
Lukislah bayangan cerminnya.
\end{exercise}

\begin{exercise}[$\star$]\label{exo:g4:symmetry:5}
Di kertas berpetak, gambarlah sumbu mendatar dan sebuah bendera
kecil (segitiga $1 \times 1$ pada tiang setinggi $3$ petak) di
atasnya. Lengkapilah gambarnya dengan simetri di bawah sumbu itu.
\end{exercise}

\begin{exercise}[$\star$]\label{exo:g4:symmetry:6}
Sebuah bangun menyentuh \omterm{def:g4:symmetry:def}{sumbu simetrinya} di titik $P$. Di mana
bayangan $P$? Jelaskan dengan gambar lipatannya.
\end{exercise}

\begin{exercise}[$\star$]\label{exo:g4:symmetry:7}
Golongkanlah huruf kapital A, B, E, F, H, N, O, T: bersumbu tegak?
bersumbu mendatar? keduanya? bukan keduanya?
\end{exercise}

\begin{exercise}[$\star$]\label{exo:g4:symmetry:8}
Sebuah segitiga punya sisi $5$ cm, $5$ cm dan $3$ cm. Apakah ia
punya \omterm{def:g4:symmetry:def}{sumbu simetri}? Sisi mana yang dipertukarkan oleh lipatan itu?
\end{exercise}

\begin{exercise}[$\star\star$]\label{exo:g4:symmetry:9}
Di kertas berpetak, gambarlah bangunmu sendiri yang punya
\emph{tepat dua} \omterm{def:g4:symmetry:def}{sumbu simetri}, lalu gambarlah kedua sumbunya.
(Petunjuk: pikirkan persegi panjang, atau tanda tambah.)
\end{exercise}

\begin{exercise}[$\star\star$]\label{exo:g4:symmetry:10}
Separuh sebuah gambar simetris diberikan: paruh kiri sebuah rumah
(dinding persegi dan separuh atap segitiga), dengan sumbu tegak
sepanjang tepi kanannya. Jelaskan atau gambarlah rumah yang lengkap.
Bagaimana panjang paruh kanan yang dilengkapi itu dibandingkan dengan paruh kirinya?
\end{exercise}

\begin{exercise}[$\star\star$]\label{exo:g4:symmetry:11}
Aya menyatakan: ``setiap segitiga dengan dua sisi sama panjang punya
\omterm{def:g4:symmetry:def}{sumbu simetri}, dan setiap segitiga yang punya \omterm{def:g4:symmetry:def}{sumbu simetri} punya
dua sisi sama panjang.'' Ujilah kedua pernyataannya pada gambar:
segitiga sama kaki, segitiga sama sisi, dan segitiga sembarang.
Apakah lipatannya mendukung dia?

\end{exercise}
